\lesson{4}{Feb 03 2022 Thu (08:18:44)}{Fascination with Fear}{Unit 3}

\subsubsection{Devices that are used by authors to create suspense}

\paragraph{Sentence Structure}

\begin{itemize}
  \item \bf{Rhythm}: a combination of long and short sentences that draw the
    reader into the experience of the story
  \item \bf{Pace}: the use of phrases, sentences, and punctuation used to build
    tension and suspense
  \item \bf{Anaphora}: the repetition of a word or phrase at the beginning of
    two or more successive clauses to emphasize an idea or create suspense
\end{itemize}

\paragraph{Time}

\begin{itemize}
  \item \bf{Foreshadowing}: the use of hints or clues that suggest future events
  \item \bf{Flashback}: an interruption in the sequence of time where a past
    event is presented in the midst of the portrayal of a current event
\end{itemize}

\paragraph{Irony}

\begin{itemize}
  \item \bf{Verbal}: a device in which what one means does not match with what
    one says
  \item \bf{Situational}: a device in which what one expects to happen does not
    actually occur
  \item \bf{Dramatic}: a device in which the audience knows something that the
    characters do not
\end{itemize}

\paragraph{Point of View}

\begin{itemize}
  \item \bf{First-person} narrators are characters in the story and participate
    in the action.
  \item \bf{Second-person} narrators address the reader as if he or she were
    part of the story's action.
  \item \bf{Third-person} omniscient narrators are outside the action but can
    give insight into the thoughts of all characters.
  \item \bf{Third-person} limited narrators are outside the action of the story
    but relate the thoughts, feelings, and experiences of one or two characters
  \item \bf{Third-person} objective narrators are outside the story; they look
    on watching every detail like a fly on the wall, but they cannot tell us
    what a character is thinking. They can only tell us what they observe.
\end{itemize}

\newpage

